%% kapa: exam=1 round=34 year=2023 period=1 form=A subject=real_estate subject_ko=부동산학원론 seq=36 number=116
%% source: Q-Net 감정평가사 기출문제 (cst00309) artlSeq=5211789 file=raw/1008/5211789/2023년도 제34회 감정평가사 1차시험 1교시.pdf
%% conversion: kordoc 4.24.51 -> tools/gen1.py (auto); review=proofread-sonnet; flags=none
\begin{question}{116}
감정평가방법 중 거래사례비교법과 관련된 설명으로 옳지 않은 것은?
\begin{choices}
  \item 거래사례비교법은 실제 거래되는 가격을 준거하므로 현실성이 있으며 설득력이 풍부하다는 장점이 있다.
  \item 거래사례비교법과 관련된 가격원칙은 대체의 원칙이고, 구해진 가액은 비준가액이라 한다.
  \item 거래사례비교법은 대상부동산과 동질·동일성이 있어서 비교 가능한 사례를 채택하는 것이 중요하다.
  \item 거래사례는 위치에 있어서 동일성 내지 유사성이 있어야 하며, 인근지역에 소재하는 경우에는 지역요인비교를 하여야 한다.
  \item 거래사례에 사정보정요인이 있는 경우 우선 사정보정을 하고, 거래시점과 기준시점간의 시간적 불일치를 정상화하는 작업인 시점수정을 하여야 한다.
\end{choices}
\end{question}
\answer{4}
